


% =============================================================================
% CITA Objective (Figure*) — NeurIPS spotlight style, math-forward + crisp
% =============================================================================
\vspace{-0.6em}
\begin{figure*}[ht!]
\centering
\begin{tcolorbox}[
  enhanced,
  colback=white,
  colframe=black,
  boxrule=0.8pt,
  borderline={0.6pt}{1.6pt}{black},
  sharp corners,
  width=0.99\textwidth,
  left=3pt,
  right=3pt,
  top=3pt,
  bottom=2pt
]
\tiny

% -------------------------
% Main objective
% -------------------------
\[
\boxed{
\mathcal{L}_{\textsc{CITA-KL}}(\theta)
\;=\;
\underbrace{
\mathbb{E}_{(I,X,Y^+,Y^-)\sim \mathcal{D}}
\Big[
-\log \sigma\!\big(\beta\,\Delta_\theta(I,X;Y^+,Y^-)\big)
\Big]
}_{\textbf{(1) Instruction-Conditioned Preference (Contrastive)}}
\;+\;
\lambda\;
\underbrace{
\mathbb{E}_{(I,X)\sim \mathcal{D}}
\Big[
\mathrm{KL}\!\big(\pi_\theta(\cdot\mid I,X)\,\|\,\pi_0(\cdot\mid I,X)\big)
\Big]
}_{\textbf{(2) Mandatory Trust-Region Anchor (KL)}}
}
\]

\vspace{-0.5em}

% -------------------------
% Definitions
% -------------------------
\[
\Delta_\theta(I,X;Y^+,Y^-)
\;=\;
\log \pi_\theta(Y^+\mid I,X)\;-\;\log \pi_\theta(Y^-\mid I,X),
\qquad
\sigma(z)=\frac{1}{1+e^{-z}}.
\]

\vspace{-0.6em}

% -------------------------
% Gradient "wow" line
% -------------------------
\[
\textbf{Gradient (preference term):}\quad
\nabla_\theta \mathcal{L}_{\textsc{pref}}
=
-\beta\,
\mathbb{E}\Big[
(1-P^+)\big(\nabla_\theta \log \pi_\theta(Y^+\!\mid I,X)-\nabla_\theta \log \pi_\theta(Y^-\!\mid I,X)\big)
\Big],
\quad
P^+=\sigma(\beta\Delta_\theta).
\]

\vspace{-0.35em}

% -------------------------
% Compact "generalized" box
% -------------------------
\[
\text{Generalized:}\qquad
\boxed{
\mathcal{L}_{\textsc{CITA-KL}}
=
\mathcal{L}_{\textsc{pref}}^{I\text{-cond}}
+
\lambda\,\mathcal{L}_{\textsc{KL}}^{I\text{-cond}},
\qquad
\lambda>0
}
\]
\end{tcolorbox}

\vspace{-0.9em}
\caption{
\textbf{CITA objective: contrastive preference learning under explicit alignment instructions.}
CITA trains on quadruples $(I,X,Y^+,Y^-)$ where the \textbf{alignment instruction $I$} \textbf{defines} the preference relation $Y^+\succ Y^-$.
\textbf{(1) Instruction-conditioned preference:} the loss is a \textbf{conditional logistic/contrastive} objective on the \textbf{score gap}
$\Delta_\theta=\log \pi_\theta(Y^+\!\mid I,X)-\log \pi_\theta(Y^-\!\mid I,X)$, with temperature $\beta$ controlling sharpness.
\textbf{(2) Mandatory KL anchor:} a \textbf{non-optional} trust-region term ($\lambda>0$) constrains updates relative to a frozen reference $\pi_0$,
stabilizing a \textbf{switchable policy family} $\{\pi_\theta(\cdot\mid I,\cdot)\}_{I\in\mathcal{I}}$ instead of collapsing to a single implicit regime.
The gradient form highlights \textbf{self-quenching updates} via $(1-P^+)$: once the instruction-conditioned preference is satisfied, the preference force diminishes, improving reliability under switching.
}
\label{fig:cita_objective_box}
\vspace{-1.2em}
\end{figure*}



% =============================================================================
% CITA: Contrastive Instruction-Tuned Alignment
% =============================================================================
% =============================================================================
% CITA: Contrastive Instruction-Tuned Alignment (math-forward, NeurIPS-spotlight)
% =============================================================================
% =============================================================================
% CITA: Contrastive Instruction-Tuned Alignment (math-forward, NeurIPS-spotlight)
% (Figure-first narrative; minimal repetition; refer to Fig.~\ref{fig:cita_objective_box})
% =============================================================================
%\vspace{-0.45em}
\section{CITA: Contrastive Instruction-Tuned Alignment}
\label{sec:cita_framework}

\textbf{Goal.} We treat alignment as \textbf{conditional control}: a single checkpoint should implement a \textbf{switchable policy family}
$\{\pi_\theta(\cdot\mid I,\cdot)\}_{I\in\mathcal{I}}$ where the \textbf{alignment instruction $I$} selects the behavioral contract at inference time.
Accordingly, CITA optimizes \textbf{instruction-indexed optima} that remain \textbf{well-conditioned under switching}.

\vspace{-0.35em}
\paragraph{Training signal: instruction makes preference identifiable.}
Each supervision unit is a preference-conditioned quadruple
\vspace{-1mm}
\[
(I, X, Y^{+}, Y^{-}) \;\in\; \mathcal{D}_{\textsc{CITA}},
\]
\vspace{-1mm}
where $Y^+\succ Y^-$ is defined \textbf{relative to $I$} for the same prompt $X$.
This explicit conditioning is the key departure from standard DPO~\cite{rafailov2023direct}: rather than learning one implicit preference regime, CITA must fit \textbf{multiple co-existing instruction-conditioned regimes}.

\vspace{-0.5em}
\paragraph{Objective (refer to Fig.~\ref{fig:cita_objective_box}).}
CITA uses the \textbf{instruction-conditioned logistic/contrastive} preference term and a \textbf{mandatory KL trust region}:
(i) the preference term increases the score gap $\Delta_\theta=\log\pi_\theta(Y^+\!\mid I,X)-\log\pi_\theta(Y^-\!\mid I,X)$ for each $(I,X)$;
(ii) the KL term anchors $\pi_\theta(\cdot\mid I,X)$ to a frozen reference $\pi_0(\cdot\mid I,X)$ with weight $\lambda>0$.
The combined effect is \textbf{instruction-specific separation} without collapsing the family into a single overconfident mode.

\vspace{-0.35em}
\paragraph{Gradient geometry: a self-quenching score-gap flow.}
Because the preference loss depends on $\theta$ only through $\Delta_\theta$, the update forms a clean vector field along the
\textbf{instruction-specific preference direction} $\nabla_\theta \Delta_\theta$:
\[
\begin{aligned}
\nabla_{\theta}\mathcal{L}_{\textsc{pref}}
=&
-\beta\,(1-P^{+})\cdot\Big(\nabla_{\theta}\log \pi_{\theta}(Y^{+}\mid I,X)\\
&-\nabla_{\theta}\log \pi_{\theta}(Y^{-}\mid I,X)
\Big)
\\[-0.15em]
P^{+}&=\sigma(\beta\Delta_\theta).
\end{aligned}
\]







This yields two switching-critical properties: \textbf{(i) self-quenching} ($P^+\!\to1 \Rightarrow \|\nabla\mathcal{L}\|\to0$), reducing over-steering across competing instructions; and
\textbf{(ii) directional purity} (a signed difference of two conditional policy gradients under the \emph{same} $(I,X)$), isolating the effect of $I$ rather than prompt artifacts. See in detail \cref{fig:cita_objective_box}.

\vspace{-0.35em}
\paragraph{Why the KL term is mandatory (trust-region view).}
Across many instructions, preference gradients can interfere, producing instruction-agnostic shortcuts or brittle oscillations.
The KL term enforces a \textbf{local trust region} around $\pi_0$:
for small parameter moves, it induces a quadratic penalty shaped by the \textbf{conditional Fisher metric},
\vspace{-1mm}
\[
\mathrm{KL}\!\big(\pi_{\theta}\,\|\,\pi_{0}\big)
\;\approx\;
\tfrac{1}{2}(\theta-\theta_0)^{\top}F_{I,X}(\theta-\theta_0),
\]
\vspace{-1mm}
so CITA performs preference optimization \textbf{on a Riemannian chart} where steps are scaled by curvature.
Practically, this keeps instruction-conditioned policies \textbf{nearby and controllable}, making switching a \textbf{stable move across neighboring optima} rather than a brittle prompt hack.



\begin{table}[ht!]
\vspace{-1em}
\centering
\scriptsize
\setlength{\tabcolsep}{4pt}
\renewcommand{\arraystretch}{1.2}
\begin{tabular}{@{}lcccc@{}}
\toprule
\textbf{Capability} & \textbf{PPO} & \textbf{GRPO} & \textbf{DPO} & \textbf{CITA} \\
\midrule
\textbf{No Reward Model} & \textcolor{red}{\ding{55}} & \textcolor{green}{\ding{51}} & \textcolor{green}{\ding{51}} & \textcolor{green}{\ding{51}} \\ \hline
\textbf{Preference Pairs} & \textcolor{red}{\ding{55}} & \textcolor{red}{\ding{55}} & \textcolor{green}{\ding{51}} & \textcolor{green}{\ding{51}} \\ \hline
\textbf{Online Generation} & \textcolor{green}{\ding{51}} & \textcolor{green}{\ding{51}} & \textcolor{red}{\ding{55}} & \textcolor{red}{\ding{55}} \\ \hline
\textbf{Instruction-Conditioned} & \textcolor{red}{\ding{55}} & \textcolor{red}{\ding{55}} & \textcolor{red}{\ding{55}} & \textcolor{green}{\ding{51}} \\ \hline
\textbf{Dynamic Policy Switch} & \textcolor{red}{\ding{55}} & \textcolor{red}{\ding{55}} & \textcolor{red}{\ding{55}} & \textcolor{green}{\ding{51}} \\ \hline
\textbf{Explicit KL Control} & \textcolor{green}{\ding{51}} & \textcolor{red}{\ding{55}} & \textbf{Implicit} & \textbf{Mandatory} \\
\bottomrule
\end{tabular}
\vspace{-1mm}
\caption{\textbf{Feature comparison across policy optimization methods.} CITA uniquely supports instruction-conditioned behavioral switching with mandatory KL regularization.}
\label{tab:feature_comparison}
\vspace{-1em}
\end{table}


\vspace{-0.35em}
\paragraph{Identifiability under prompt-held-constant switching.}
ECLIPTICA-style construction creates \emph{paired} conditions $(I_a,X)$ and $(I_b,X)$ with incompatible behavioral targets.
Any instruction-invariant policy $\pi(\cdot\mid X)$ cannot simultaneously satisfy both preference relations, incurring unavoidable loss.
Thus, the only degree of freedom that can explain systematic behavioral changes is the \textbf{control variable $I$}, turning alignment into a \textbf{causal control problem}.


\begin{comment}
    
\paragraph{Hyperparameter optimization.}
We tune CITA with \textbf{Optuna}~\cite{akiba2019optuna} using the \textbf{TPE} sampler for \textbf{12 trials}, following recent HPO guidance~\cite{liang2023holistic}.
A consistent trend emerges: \textbf{CITA\_Instruct prefers a $\sim$50\% lower learning rate} than CITA\_NoInstruct.
Mechanistically, instruction-augmented inputs are \textbf{30--40\% longer}, which increases effective gradient magnitude; a smaller step size stabilizes the preference/KL trade-off.

\vspace{-0.25em}
\begin{table}[t]
\centering
\footnotesize
\setlength{\tabcolsep}{6pt}
\renewcommand{\arraystretch}{1.06}
\resizebox{0.98\columnwidth}{!}{%
\begin{tabular}{@{}lcc@{}}
\toprule
\textbf{Hyperparameter} & \textbf{CITA\_NoInstruct (T5)} & \textbf{CITA\_Instruct (T7)} \\
\midrule
Learning rate          & 6.83e-6  & 5.41e-6  \\
$\lambda_{\text{KL}}$  & 0.00052  & 0.00023  \\
$\beta$ (temp.)        & 0.119    & 0.107    \\
Weight decay           & 0.0091   & 0.0109   \\
Warmup ratio           & 7.5\%    & 10.0\%   \\
\midrule
Final margin           & 6.95     & 7.52     \\
Accuracy               & 89.5\%   & 89.0\%   \\
\bottomrule
\end{tabular}
}
\vspace{-0.35em}
\caption{\textbf{Optuna-selected hyperparameters.} T5/T7 denote the best-performing trial IDs for each setting.}
\label{tab:hyperparams}
\vspace{-1.5em}
\end{table}
\end{comment}




\vspace{-0.35em}
\paragraph{Takeaway.}
\textbf{CITA optimizes an instruction-indexed family of nearby optima, with the KL anchor supplying the geometry that makes switching reliable; Table~\ref{tab:feature_comparison} summarizes how this differs from prior alignment methods.}


















\begin{comment}

\begin{figure}[H]
\centering
\small
\setlength{\tabcolsep}{2pt}
\begin{tabular}{ccccccc}
\textit{\scriptsize Llama-3.1-8B} & & \textit{\scriptsize PKU (chosen)} & & \textit{\scriptsize PKU (pairs)} & & \textit{\scriptsize + Instructions} \\[2pt]
$\downarrow$ & & $\downarrow$ & & $\downarrow$ & & $\downarrow$ \\[2pt]
\fbox{\textbf{Base}} & $\rightarrow$ & \fbox{\textbf{SFT}} & $\rightarrow$ & \fbox{\textbf{DPO}} & $\rightarrow$ & \colorbox{green!20}{\fbox{\textbf{CITA}}} \\[2pt]
& & $\downarrow$ & & $\downarrow$ & & $\downarrow$ \\[2pt]
& & $\mathcal{L}_{\text{SFT}}$ & & $\mathcal{L}_{\text{DPO}}$ & & $\mathcal{L}_{\text{CITA}} + \lambda \cdot \text{KL}$ \\
\end{tabular}
\caption{CITA training pipeline. Each stage uses the previous checkpoint. CITA adds mandatory KL regularization for instruction-conditioned alignment.}
\label{fig:pipeline}
\end{figure}

\end{comment}




